\chapter{Metodología y diseño experimental}
\label{ch:metodologia}
% Núcleo pre-registrado (redactar antes de ver resultados). Fuente: 1 - Diseño.md + pending/.
% Frontera con implementación: aquí el QUÉ y el POR QUÉ (decisiones de diseño);
% el CÓMO (realización en software) va en el cap. \ref{ch:implementacion}.

Este capítulo traduce el posicionamiento del Capítulo~\ref{ch:estado-del-arte} en un diseño experimental contrastable. La Sección~\ref{sec:hipotesis} formula las seis hipótesis, y la Sección~\ref{sec:variables} define las variables que las operacionalizan. La Sección~\ref{sec:ventana} fija la ventana temporal de medición, y las Secciones~\ref{sec:setup} y~\ref{sec:matriz} describen la configuración del entrenamiento y la matriz de experimentos. La Sección~\ref{sec:calibracion} detalla cómo se calibraron los presupuestos de \emph{epochs} y los umbrales de \emph{accuracy} mediante un piloto previo a la matriz. Cierran el capítulo el protocolo de evaluación (Sección~\ref{sec:protocolo-eval}), el protocolo de análisis (Sección~\ref{sec:protocolo-analisis}) y los riesgos del diseño (Sección~\ref{sec:riesgos}).

\section{Hipótesis}
\label{sec:hipotesis}

% Pendiente: las seis hipótesis se reescriben sin criterio de decisión cuando se defina el método de análisis.

\section{Variables}
\label{sec:variables}

\subsection{Dependientes: eficiencia del entrenamiento}
% Tres constructos. Velocidad: VD1 epochs-a-umbral sobre val suavizada (primaria, censurado),
% VD2 AUC de val-loss, VD3 mejor val-loss. Rendimiento final: VD4 final_test_acc (+ F1-macro).
% Generalización: VD5 gap_loss = final_test_loss - final_train_eval_loss (primaria),
% VD6 gap_acc = final_train_eval_acc - final_test_acc (robustez). Detalle en 1 - Diseño.md.

\subsection{Independientes: métricas tempranas}
% Las dos familias; lista cerrada antes de ejecutar (anti p-hacking).

\section{Ventana temporal}
\label{sec:ventana}
% Fracciones 5/10/25/50 %; el barrido es resultado reportable.

\section{Configuración del entrenamiento}
\label{sec:setup}
% Pendiente en esta sección: arquitecturas (FC/CNN/ResNet-18), optimizadores (SGD/Adam) y
% normalización sin augmentation.

El estudio emplea cuatro conjuntos de clasificación de imágenes de dificultad creciente: MNIST~\cite{lecun1998gradient}, CIFAR-10 y CIFAR-100~\cite{krizhevsky2009learning} y Tiny-ImageNet~\cite{le2015tiny}. En todos ellos el conjunto de test oficial se mantiene intacto y el de validación se recorta del entrenamiento oficial mediante muestreo estratificado por clase con una \emph{seed} fija; en CIFAR esto reproduce la partición 45.000/5.000 de He et al.~\cite{he2015deep}. Así la validación conserva las proporciones de clase del entrenamiento y la partición es idéntica en los 960 entrenamientos. La Tabla~\ref{tab:datasets} recoge los tamaños resultantes, el número de clases y la dimensión de entrada de cada conjunto. Tiny-ImageNet es la excepción: como las etiquetas de su conjunto de test oficial no son públicas, su conjunto de validación público se emplea como test, y por eso el total usable baja a 110.000 ejemplos en lugar de 120.000.

\begin{table}[ht]
\centering
\small
\caption{Conjuntos de datos y particiones.}
\label{tab:datasets}
\begin{tabular}{lrrrrcc}
\toprule
\textbf{Conjunto} & \textbf{Total} & \textbf{Entren.} & \textbf{Valid.} & \textbf{Test} & \textbf{Clases} & \textbf{Dimensión} \\
\midrule
MNIST          & 70.000  & 50.000 & 10.000 & 10.000 & 10  & $1{\times}28{\times}28$ \\
CIFAR-10       & 60.000  & 45.000 & 5.000  & 10.000 & 10  & $3{\times}32{\times}32$ \\
CIFAR-100      & 60.000  & 45.000 & 5.000  & 10.000 & 100 & $3{\times}32{\times}32$ \\
Tiny-ImageNet  & 110.000 & 90.000 & 10.000 & 10.000 & 200 & $3{\times}64{\times}64$ \\
\bottomrule
\end{tabular}
\end{table}

La Tabla~\ref{tab:datasets} fija el tamaño del problema; conviene además situar qué \emph{accuracy} cabe esperar de cada arquitectura en estas condiciones. La comparación se hace contra el \emph{régimen comparable} (redes entrenadas desde cero, sin \emph{data augmentation} y a resolución nativa), que es el de este estudio. No se compara contra el estado del arte, que casi siempre usa \emph{data augmentation}, \emph{weight decay} y entrenamientos largos. La Tabla~\ref{tab:referencias} reúne ese régimen para las tres arquitecturas, con cifras tomadas de filas de ablación o de referencia de la literatura: para la red totalmente conectada (un perceptrón multicapa, MLP), Bachmann et al.~\cite{bachmann2023scaling}; para la CNN pequeña, Chen y Ho~\cite{chen2018deep} (cota superior, con activación GReLU) y Mantzaris~\cite{mantzaris2025exploring} (cota inferior, con \emph{batch normalization}); y para ResNet-18, las filas sin \emph{data augmentation} de Liu et al.~\cite{liu2024zeropur} y Ko et al.~\cite{ko2024learning}, con Larsson et al.~\cite{larsson2017fractalnet} como corroboración en CIFAR-100. En MNIST cualquier red pequeña satura en torno al 99\,\% (LeCun et al.~\cite{lecun1998gradient}).

Esas cifras se leen como bandas esperadas y no como objetivos, porque las fuentes citadas suelen entrenar durante más tiempo o con redes algo mayores y por tanto acotan por arriba. Las dos entradas marcadas con $\dagger$ corresponden a la CNN simple y a la ResNet-18 en Tiny-ImageNet, para las que no existe referencia publicada en este régimen (desde cero, sin \emph{data augmentation} y a resolución nativa de 64$\times$64); son estimaciones propias, y los resultados de este estudio aportarán las primeras cifras en ese régimen.

\begin{table}[ht]
\centering
\small
\caption{Accuracy de test de referencia por arquitectura.}
\label{tab:referencias}
\begin{tabular}{lcccc}
\toprule
\textbf{Arquitectura} & \textbf{MNIST} & \textbf{CIFAR-10} & \textbf{CIFAR-100} & \textbf{Tiny-ImageNet} \\
\midrule
FC (MLP 1024-1024)          & $\sim$98,4\,\%       & $\sim$50--54\,\% & $\sim$25--29\,\% & $\sim$8--12\,\% \\
CNN (3 bloques)             & $\sim$99,0--99,4\,\% & $\sim$70--80\,\% & $\sim$45--55\,\% & $\sim$20--35\,\%\,$\dagger$ \\
ResNet-18 (stem 3$\times$3) & $\sim$99,0--99,5\,\% & $\sim$84--92\,\% & $\sim$50--57\,\% & $\sim$40--50\,\%\,$\dagger$ \\
\bottomrule
\end{tabular}
\end{table}

Como cota laxa de contexto, ResNet-18 entrenada desde cero a resolución nativa pero \emph{con} \emph{data augmentation} alcanza 94,55\,\% (CIFAR-10), 78,07\,\% (CIFAR-100) y 50,68\,\% (Tiny-ImageNet) según GenURL~\cite{li2023genurl}; la distancia con la Tabla~\ref{tab:referencias} es justo lo que aportan el \emph{data augmentation}, el \emph{weight decay} y los entrenamientos largos. Para el MLP y la CNN, los mejores resultados publicados exigen cambiar de clase de arquitectura (redes densas que aprenden conectividad local tipo convolución, Neyshabur~\cite{neyshabur2020towards}; redes totalmente convolucionales profundas, Springenberg et al.~\cite{springenberg2015striving}), de modo que no son techos comparables del MLP llano ni de la CNN simple.

\section{Matriz de experimentos}
\label{sec:matriz}
% Fuente única de la rejilla: src/config.py (DATASETS, MODELS, OPTIMIZERS, SEEDS, LR_GRID,
% FIXED_KNOBS). Toda cifra de esta sección debe cuadrar con esas constantes.

La matriz de experimentos determina qué entrenamientos se ejecutan y qué combinaciones de hiperparámetros se barren. Combina tres factores, descritos en la Sección~\ref{sec:setup}: cuatro conjuntos de datos, tres arquitecturas y dos optimizadores. Cada combinación de los tres forma una celda, y hay 24 en total. Además, dentro de cada celda se barren ocho \emph{learning rates} y cinco \emph{seeds}, lo que da 40 entrenamientos por celda y 960 en total. La Tabla~\ref{tab:matriz} recoge los valores de cada eje.

La celda es la unidad de análisis: dentro de ella el conjunto de datos, la arquitectura y el optimizador permanecen fijos, y lo único que cambia entre sus 40 entrenamientos es el \emph{learning rate} y la \emph{seed}. Las correlaciones se calculan dentro de cada celda y no sobre los 960 entrenamientos juntos. Mezclarlos haría que casi toda la variación en eficiencia procediese del conjunto de datos y de la arquitectura, porque todo entrenamiento sobre MNIST es rápido y ninguno sobre Tiny-ImageNet lo es, y la correlación resultante sería alta aunque la métrica no distinguiese nada dentro de un problema concreto. Con los tres factores fijos, en cambio, lo único que separa en eficiencia a dos entrenamientos de la misma celda es su \emph{learning rate}, y la correlación mide entonces lo que el trabajo pregunta: si la métrica temprana anticipa qué \emph{learning rates} van a entrenar bien. El procedimiento se detalla en la Sección~\ref{sec:protocolo-analisis}.

\begin{table}[ht]
\centering
\small
\setlength{\tabcolsep}{4pt}
\caption{Ejes de la matriz de experimentos.}
\label{tab:matriz}
\begin{tabular}{llc}
\toprule
\textbf{Eje} & \textbf{Valores} & \textbf{Niveles} \\
\midrule
Conjunto de datos & MNIST, CIFAR-10, CIFAR-100, Tiny-ImageNet & 4 \\
Arquitectura      & FC, CNN, ResNet-18 & 3 \\
Optimizador       & SGD, Adam & 2 \\
\addlinespace
\emph{Learning rate} (SGD)  & $3 \times 10^{-4}$, $10^{-3}$, $3 \times 10^{-3}$, $10^{-2}$, $3 \times 10^{-2}$, $10^{-1}$, $3 \times 10^{-1}$, $1$ & 8 \\
\emph{Learning rate} (Adam) & $3 \times 10^{-5}$, $10^{-4}$, $3 \times 10^{-4}$, $10^{-3}$, $3 \times 10^{-3}$, $10^{-2}$, $3 \times 10^{-2}$, $10^{-1}$ & 8 \\
\emph{Seed}                 & 0, 1, 2, 3, 4 & 5 \\
\bottomrule
\end{tabular}
\end{table}

Los ocho \emph{learning rates} están espaciados de forma logarítmica, multiplicando cada valor por algo más de tres para obtener el siguiente, porque lo que separa a un \emph{learning rate} de otro es su orden de magnitud y no la diferencia entre ambos. El recorrido es ancho a propósito, con el valor mayor unas 3.000 veces por encima del menor, de manera que la lista abarca desde entrenamientos demasiado lentos hasta entrenamientos que divergen. Esos extremos no son cómputo desperdiciado, sino lo que da rango a la variable que se quiere predecir: si todas las configuraciones fuesen razonables, la eficiencia sería casi constante y no habría forma de distinguir una métrica predictiva de una que no lo es. Los entrenamientos que no alcanzan el umbral quedan censurados en la velocidad de convergencia, en el sentido de la Sección~\ref{sec:eficiencia}.

Cada optimizador tiene su propia lista, y las tres arquitecturas comparten la del optimizador con el que entrenan. Dar una lista distinta a cada arquitectura, centrada donde mejor le funcione, exigiría conocer ese óptimo antes de ejecutar los entrenamientos, es decir, fijar el diseño con información extraída del resultado. La lista de Adam es la de SGD dividida entre 10, porque en Adam el paso efectivo mide aproximadamente el \emph{learning rate}, como expone la Sección~\ref{sec:sgd}, mientras que en SGD el \emph{momentum} de 0,9 lo amplifica alrededor de 10 veces; un mismo valor nominal resulta así un paso corriente en SGD y un salto desmedido en Adam.

Salvo por dónde empiezan, las dos listas son idénticas en número de valores, en factor entre valores consecutivos y en anchura de recorrido. Esa simetría permite emparejar SGD y Adam por posición en la lista, que es lo que sitúa a ambos en el mismo tramo de su propio rango útil, y es lo que exige el objetivo OE5. Emparejarlos por valor nominal parecería más limpio, pero compararía tramos distintos, porque un mismo número deja a SGD todavía trabajando y a Adam muy por encima de su rango.

Las cinco \emph{seeds} son las mismas en todas las celdas. Una \emph{seed} fija la inicialización de los pesos, el orden de los \emph{minibatches} y qué ejemplos forman el \emph{batch} de medición, pero no cambia el régimen del entrenamiento, cosa que el \emph{learning rate} sí hace. Por eso el barrido se reparte a favor de los \emph{learning rates} y las \emph{seeds} se limitan a las necesarias para medir la dispersión que queda al fijar todo lo demás. Compartirlas entre celdas permite además comparar SGD y Adam con idéntica inicialización e idéntico orden de los datos.

El resto de hiperparámetros toma el mismo valor en los 960 entrenamientos: tamaño de \emph{batch} 128, \emph{momentum} 0,9 en SGD, \emph{weight decay} nulo y $M = 256$ ejemplos en el \emph{batch} de medición. Al no variar ninguno, ninguno puede actuar como variable de confusión. La excepción son el presupuesto de \emph{epochs} y el umbral de \emph{accuracy}, que dependen del conjunto de datos y por tanto de un eje de la matriz; su calibración se detalla en la Sección~\ref{sec:calibracion}.

\section{Calibración de presupuestos y umbrales}
\label{sec:calibracion}
% Piloto de calibración: protocolo, criterios preescritos y valores finales (2026-06-17).
% Fuente: 2 - Decisiones.md (2026-06-09) + 3 - Progreso.md. La realización software del piloto
% va en el cap. \ref{ch:implementacion} (\ref{sec:pilot}).

El presupuesto de \emph{epochs} $T$ y el umbral de \emph{accuracy} de validación $\tau$ de cada conjunto de datos se fijan en dos pasos: se parte de valores candidatos y un piloto de calibración, ejecutado antes de lanzar la matriz, los confirma o corrige sobre curvas de entrenamiento reales. El piloto consta de un entrenamiento por celda de la matriz (cada combinación de conjunto de datos, arquitectura y optimizador de la Sección~\ref{sec:matriz}), es decir, 24 en total, con el \emph{learning rate} en el centro de la rejilla de cada optimizador ($10^{-2}$ para SGD y $10^{-3}$ para Adam, los valores por defecto canónicos), \emph{seed} 0 y el doble del presupuesto candidato de \emph{epochs} (40, 80, 120 o 160 según el conjunto de datos). La duplicación responde a una asimetría de costes: recortar a posteriori una curva generosa no cuesta nada, mientras que alargar una curva corta obliga a relanzar el entrenamiento.

Para evitar circularidad, los criterios de decisión se dejaron escritos antes de ejecutar el piloto. El presupuesto final se fija donde el \emph{loss} de validación alcanza su meseta, más un margen, redondeado a un múltiplo de 20 para que las fracciones de ventana de la Sección~\ref{sec:ventana} (5, 10, 25, 50 y 100\,\%) caigan exactamente en frontera de \emph{epoch}. El umbral de \emph{accuracy} se fija de modo que la CNN y la ResNet-18 con el \emph{learning rate} centrado lo crucen hacia el 30--60\,\% del presupuesto, es decir, que sea alcanzable pero no trivial: un umbral que se cruza en la primera \emph{epoch} no discrimina velocidad, y uno que casi ningún entrenamiento cruza censura media matriz. La Tabla~\ref{tab:calibracion} recoge, para cada conjunto de datos, el presupuesto de \emph{epochs} $T$ y el umbral de \emph{accuracy} de validación $\tau$, tanto en su valor candidato previo al piloto como en el finalmente adoptado tras él.

\begin{table}[ht]
\centering
\small
\caption{Calibración por conjunto de datos.}
\label{tab:calibracion}
\begin{tabular}{lcccc}
\toprule
& \multicolumn{2}{c}{\textbf{Presupuesto $T$ (\emph{epochs})}} & \multicolumn{2}{c}{\textbf{Umbral $\tau$}} \\
\cmidrule(lr){2-3}\cmidrule(lr){4-5}
\textbf{Conjunto} & Candidato & Final & Candidato & Final \\
\midrule
MNIST          & 20 & 20 & 0,97 & 0,97 \\
CIFAR-10       & 40 & 40 & 0,75 & 0,65 \\
CIFAR-100      & 60 & 40 & 0,35 & 0,35 \\
Tiny-ImageNet  & 80 & 40 & 0,25 & 0,20 \\
\bottomrule
\end{tabular}
\end{table}

Dos umbrales bajaron respecto a su candidato. En CIFAR-10 el umbral pasó de 0,75 a 0,65: el techo de \emph{accuracy} de validación de la CNN con el \emph{learning rate} centrado ronda 0,72--0,73, de modo que con 0,75 esa celda quedaba censurada de forma permanente, no por lentitud sino por techo. En Tiny-ImageNet pasó de 0,25 a 0,20: el margen de la CNN sobre 0,25 era de apenas 0,007, frágil frente a la varianza entre \emph{seeds}, mientras que con 0,20 sube a en torno a 0,06. Dos presupuestos se recortaron, el de CIFAR-100 de 60 a 40 \emph{epochs} y el de Tiny-ImageNet de 80 a 40: sin \emph{data augmentation} las curvas de \emph{loss} de validación se aplanan mucho antes del presupuesto candidato, y cada \emph{epoch} sobrante sería coste muerto multiplicado por los 960 entrenamientos de la matriz.

La calibración no altera el protocolo de medición: todas las métricas se calculan al final de cada \emph{epoch} sobre el \emph{batch} de medición fijo, en todos los entrenamientos, y conservar la serie temporal completa permite escoger las ventanas de análisis a posteriori.

\section{Protocolo de evaluación}
\label{sec:protocolo-eval}
% Split train/val/test (confirmado por tutor 2026-06-12).
% Gap de generalización como variable objetivo (confirmado por tutor 2026-06-14, implementado en src/).

\section{Protocolo de análisis}
\label{sec:protocolo-analisis}
% Pendiente: se reconstruye objetivo por objetivo (fase B).
% Viene de \ref{sec:matriz}, que remite aquí y por tanto DEBE quedar cubierto: cada uno de los 40
% entrenamientos de una celda aporta un par de cifras, el valor de la métrica en sus primeras
% epochs y el indicador de eficiencia del entrenamiento completo; la correlación de la celda se
% calcula entre esas dos series de 40 valores; el cálculo se repite en las 24 celdas y la pregunta
% pasa a ser si todas dan la misma respuesta.

\section{Riesgos}
\label{sec:riesgos}
% Pendiente: paradoja de Simpson al agregar.

% Plantilla de capítulo: cerrar con \section{Conclusiones del capítulo} cuando el cuerpo esté
% escrito, al modo de \ref{sec:fund-conclusiones}. No crear el título antes que la prosa.
